%-----------------------------------------------------------------------------
%  Camera Power Adapter - Custom Datasheet & Design Calculations
%  Document: DesignDocs/CameraPowerAdapter_DesignDocument.tex
%
%  Compile:  pdflatex CameraPowerAdapter_DesignDocument.tex  (run twice)
%         or latexmk -pdf CameraPowerAdapter_DesignDocument.tex
%
%  Template notes:
%   - \todo{...} marks content that still needs to be filled in / verified.
%   - Numeric examples carried over from BuckBoostTopology.md and
%     ComponentSelection.md are marked "provisional" until re-verified
%     against the finalized schematic.
%-----------------------------------------------------------------------------
\documentclass[11pt,a4paper]{article}

%--------------------------------- Packages ---------------------------------
\usepackage[utf8]{inputenc}
\usepackage[T1]{fontenc}
\usepackage{lmodern}
\usepackage[margin=22mm,headheight=15pt]{geometry}
\usepackage{amsmath}
\usepackage{amssymb}
\usepackage{siunitx}
\usepackage{booktabs}
\usepackage{longtable}
\usepackage{array}
\usepackage{multirow}
\usepackage{tabularx}
\usepackage{float}
\usepackage{graphicx}
\usepackage{xcolor}
\usepackage{fancyhdr}
\usepackage{enumitem}
\usepackage[colorlinks=true,
            linkcolor=blue!60!black,
            urlcolor=blue!60!black,
            citecolor=blue!60!black]{hyperref}

\sisetup{per-mode=symbol}
\setlength{\emergencystretch}{2em}   % avoid overfull lines from long \texttt words

%------------------------------ Custom macros -------------------------------
\newcommand{\todo}[1]{\textcolor{red}{\textbf{[TODO: #1]}}}
\newcommand{\prov}{\textcolor{orange!80!black}{\textbf{(provisional)}}}
\newcommand{\note}{\par\medskip\noindent\textbf{Note ---}\ }
\newcommand{\comp}[1]{\texttt{#1}}                 % reference designator
\newcommand{\figplaceholder}[1]{%
  \fbox{\parbox[c][45mm][c]{0.85\linewidth}{\centering\color{gray}#1}}}

%------------------------------ Header / footer -----------------------------
\pagestyle{fancy}
\fancyhf{}
\fancyhead[L]{Camera Power Adapter}
\fancyhead[R]{Custom Datasheet \& Design Calculations}
\fancyfoot[C]{\thepage}
\renewcommand{\headrulewidth}{0.4pt}

%=============================================================================
\begin{document}

%------------------------------- Title page ---------------------------------
\begin{titlepage}
  \hypersetup{pageanchor=false}
  \centering
  \vspace*{2cm}

  {\Large Vintage Film Camera Equipment\par}
  \vspace{0.5cm}
  {\Huge\bfseries Dual-Input Priority-Switched\\[0.3cm] Buck-Boost Power Adapter\par}
  \vspace{1.0cm}
  {\LARGE Custom Datasheet \& Design Calculations\par}
  \vspace{2.0cm}

  \begin{tabularx}{0.98\textwidth}{@{}lX@{}}
    \toprule
    \textbf{Document ID}     & CPA-DD-001 \\[2pt]
    \textbf{Revision}        & 0.1 --- Draft / Template \\[2pt]
    \textbf{Date}            & \today \\[2pt]
    \textbf{Author}          & \todo{name} \\[2pt]
    \textbf{Status}          & Work in progress --- sections marked \todo{} are incomplete \\[2pt]
    \textbf{Related spec}    & \texttt{Specifications/SpecSheet.pdf} \\[2pt]
    \textbf{Supporting docs} & \texttt{BuckBoostTopology.md}, \texttt{ComponentSelection.md}, \texttt{SchematicBOM.md} \\
    \bottomrule
  \end{tabularx}

  \vfill
  {\small This document defines the board as a standalone power module (custom datasheet) and
   records the calculations behind every selected value (design calculations).}
\end{titlepage}
\hypersetup{pageanchor=true}

\tableofcontents
\newpage

%=============================================================================
\section{Introduction}
%=============================================================================

\subsection{Purpose}
This document serves two functions:
\begin{enumerate}[nosep]
  \item A \textbf{custom datasheet} for the assembled adapter --- what the board
        does, its limits, and how to connect it.
  \item A \textbf{design calculation record} --- every value on the schematic is
        derived from an equation traceable to a manufacturer datasheet.
\end{enumerate}

\subsection{Scope}
The adapter accepts two independent \SIrange{11}{36}{\V} DC inputs with fixed
priority (Input~2 preferred), protects against reverse polarity and overcurrent,
and delivers a jumper-selected \SI{12}{\V} or \SI{24}{\V} output rated
\SIrange{6}{8}{\A} (burst duty cycle; see \autoref{sec:datasheet-operating}).
\newline
\todo{Confirm scope wording against the final schematic.}

\subsection{Related documents}
\begin{table}[H]
  \centering
  \begin{tabularx}{\linewidth}{@{}lX@{}}
    \toprule
    \textbf{Document} & \textbf{Contents} \\
    \midrule
    \texttt{Specifications/SpecSheet.pdf}   & Customer requirements / acceptance criteria \\
    \texttt{BuckBoostTopology.md}           & LT8390 stage build guide and audit log \\
    \texttt{ComponentSelection.md}          & Value-by-value selection rationale \\
    \texttt{SchematicBOM.md}                & Current bill of materials and placement status \\
    \texttt{Datasheets/}                    & Manufacturer datasheets (local copies) \\
    \bottomrule
  \end{tabularx}
  \caption{Reference documents.}
  \label{tab:related-docs}
\end{table}

\subsection{Acronyms}
\begin{table}[H]
  \centering
  \begin{tabularx}{\linewidth}{@{}lX@{}}
    \toprule
    \textbf{Term} & \textbf{Meaning} \\
    \midrule
    UVLO    & Under-voltage lockout \\
    CC      & Constant-current (limit) \\
    HS / LS & High-side / low-side MOSFET switch \\
    PGOOD   & Power-good output of the LT8390 (active low) \\
    NTC     & Negative temperature coefficient thermistor \\
    \bottomrule
  \end{tabularx}
  \caption{Acronyms and definitions.}
\end{table}

%=============================================================================
\section{System Overview}
%=============================================================================

\subsection{Block diagram}
\begin{figure}[H]
  \centering
  \figplaceholder{\textbf{Figure placeholder}\\[2mm]
    Input 1 $\rightarrow$ [fuse, reverse-pol.] $\;$\\
    Input 2 $\rightarrow$ [fuse, reverse-pol.] $\;$
    $\bigg\}$ Priority mux (LTC4417) $\rightarrow$ LT8390 buck-boost
    $\rightarrow$ 12\,V/24\,V output\\
    \small Export a block diagram from KiCad or draw with TikZ and \texttt{\textbackslash includegraphics} here.}
  \caption{System block diagram. \todo{replace placeholder.}}
  \label{fig:blockdiagram}
\end{figure}

\subsection{Operating modes}
\begin{table}[H]
  \centering
  \begin{tabularx}{\linewidth}{@{}lX@{}}
    \toprule
    \textbf{Mode} & \textbf{Behaviour} \\
    \midrule
    Input priority   & Input~2 powers the board whenever it is valid; automatic fallback to Input~1. \\
    Buck region      & $V_{IN} > V_{OUT}$ --- buck leg switches. \\
    Boost region     & $V_{IN} < V_{OUT}$ --- boost leg switches. \\
    Overlap band     & $V_{IN} \approx V_{OUT}$ --- both legs switch (worst thermal corner). \\
    Output select    & Jumper fitted = \SI{12}{\V}; jumper removed = \SI{24}{\V}. \\
    Over-temp fault  & Latch pulls EN/UVLO low until input power is cycled. \\
    \bottomrule
  \end{tabularx}
  \caption{Operating modes.}
\end{table}

%=============================================================================
\section{Custom Datasheet}
\label{sec:datasheet}
%=============================================================================

\subsection{Description}
Dual-input, priority-selected, non-isolated synchronous buck-boost DC/DC adapter
for vintage film camera equipment. Intended for low-duty-cycle loads (short
bursts) inside a small enclosure --- the enclosure floor acts as the heat sink
through a bottom-side thermal pad.

\subsection{Absolute maximum ratings}
\begin{table}[H]
  \centering
  \begin{tabularx}{\linewidth}{@{}lXl@{}}
    \toprule
    \textbf{Parameter} & \textbf{Conditions / note} & \textbf{Max} \\
    \midrule
    Input voltage, either input   & Continuous; reverse-polarity block extends to $-42$\,V      & \SI{36}{\V} \\
    Output current                & Burst rating                                                & \SI{8}{\A} \\
    Output current                & Continuous duty (thermal-limited)                           & \todo{\SI{2}{\A}?} \\
    Operating ambient temperature & \todo{state enclosure assumption}                           & \todo{} \\
    Storage temperature           &                                                             & \todo{} \\
    \bottomrule
  \end{tabularx}
  \caption{Absolute maximum ratings. Stresses beyond these may damage the board.}
\end{table}

\subsection{Recommended operating conditions}
\label{sec:datasheet-operating}
\begin{table}[H]
  \centering
  \begin{tabularx}{\linewidth}{@{}lXl@{}}
    \toprule
    \textbf{Parameter} & \textbf{Conditions / note} & \textbf{Value} \\
    \midrule
    Input 1 voltage & Valid operating window (UVLO $\approx$ \SI{11}{\V}) & \SIrange{11}{36}{\V} \\
    Input 2 voltage & Valid operating window (UVLO $\approx$ \SI{11}{\V}) & \SIrange{11}{36}{\V} \\
    Output voltage  & Jumper selection                                            & \SI{12}{\V} or \SI{24}{\V} \\
    Output current  & Real duty cycle (per spec)                                  & \SIrange{2}{4}{\A} typical \\
    Output current  & Short-burst rating                                          & \SIrange{6}{8}{\A} \\
    Ambient         & Inside camera enclosure                                      & \todo{} \\
    \bottomrule
  \end{tabularx}
  \caption{Recommended operating conditions.}
\end{table}

\subsection{Electrical characteristics}
\todo{fill min/typ/max columns once the first board is measured; values below are
design targets, not tested limits.}
\begin{table}[H]
  \centering
  \begin{tabularx}{\linewidth}{@{}lXl@{}}
    \toprule
    \textbf{Parameter} & \textbf{Conditions / note} & \textbf{Target} \\
    \midrule
    \multicolumn{3}{@{}l}{\textbf{Inputs}}\\
    \addlinespace[1pt]
    Input priority            & Fixed, not voltage-based                          & Input~2 $>$ Input~1 \\
    UVLO per channel          & Set by priority controller; below mux floor       & \SI{11}{\V} $\pm$ \todo{} \\
    Reverse polarity          & Both inputs independently                         & \SI{-42}{\V} capable \\
    Fuse per input line       & Sized from worst-case input current               & \todo{rating} \\
    Input current (worst case)& $P_{OUT}/(\eta \cdot V_{IN,min})$                  & \todo{value in A} \\
    \addlinespace[3pt]
    \multicolumn{3}{@{}l}{\textbf{Output}}\\
    \addlinespace[1pt]
    Output voltage accuracy   & LT8390 FB accuracy                                & $\pm$1.5\,\% \\
    Output current limit      & $R_{11} = \SI{12}{\m\ohm}$ \prov                    & \SI{8.33}{\A} \\
    Output ripple             & Measured at rated burst load                      & \todo{\si{mV_{pp}}} \\
    \addlinespace[3pt]
    \multicolumn{3}{@{}l}{\textbf{Protection \& indicators}}\\
    \addlinespace[1pt]
    Over-temperature trip     & NTC/TL431 latch \prov                              & $\approx$\SI{85}{\celsius} $\pm$ \SI{8}{\celsius} \\
    Fault recovery            & Latch-off; requires input power cycle             & --- \\
    Indicators                & Green = output OK, red = fault, red = over-temp   & see \autoref{tab:leds} \\
    \bottomrule
  \end{tabularx}
  \caption{Electrical characteristics (design targets).}
\end{table}

\subsection{Connector pinout}
\begin{table}[H]
  \centering
  \begin{tabular}{@{}cll@{}}
    \toprule
    \textbf{Designator} & \textbf{Signal} & \textbf{Note} \\
    \midrule
    IN1 & Input 1 positive & M3 solder-on threaded standoff \\
    IN1 & Input 1 return   & M3 solder-on threaded standoff \\
    IN2 & Input 2 positive & M3 solder-on threaded standoff \\
    IN2 & Input 2 return   & M3 solder-on threaded standoff \\
    OUT & Output positive (V+) & M3 solder-on threaded standoff \\
    OUT & Output return (GND)  & M3 solder-on threaded standoff \\
    \bottomrule
  \end{tabular}
  \caption{I/O connections. \todo{confirm silkscreen labels once the layout is done.}}
\end{table}

\subsection{Mechanical \& environmental}
\begin{table}[H]
  \centering
  \begin{tabularx}{\linewidth}{@{}lX@{}}
    \toprule
    \textbf{Parameter} & \textbf{Value / requirement} \\
    \midrule
    PCB outline            & \SI{45}{\mm} $\times$ \SI{75}{\mm} max \\
    Mounting holes         & 4 $\times$ \SI{3}{\mm} dia., \SI{70}{\mm} $\times$ \SI{40}{\mm} pattern \\
    Bottom thermal pad     & \SI{45}{\mm} $\times$ \SI{40}{\mm}, contact with enclosure floor \\
    Assembled height       & $<\SI{13.5}{\mm}$ target; \SI{18}{\mm} fallback ceiling \\
    Indicators             & Bottom side, output end \\
    \bottomrule
  \end{tabularx}
  \caption{Mechanical requirements.}
\end{table}

\subsection{Indicator logic}
\begin{table}[H]
  \centering
  \begin{tabular}{@{}lll@{}}
    \toprule
    \textbf{Green (power)} & \textbf{Red (fault)} & \textbf{Red (over-temp)} \\
    \midrule
    ON  & OFF  & OFF \\
    OFF & ON   & OFF \\
    OFF & OFF  & ON  \\
    \bottomrule
  \end{tabular}
  \caption{Indicator truth table. \todo{add state column once LED wiring is final.}}
  \label{tab:leds}
\end{table}

\subsection{Typical application}
\begin{figure}[H]
  \centering
  \figplaceholder{Place the annotated top-level schematic or KiCad plot here.}
  \caption{Typical application circuit. \todo{replace placeholder.}}
\end{figure}

%=============================================================================
\section{Design Calculations}
\label{sec:calcs}
%=============================================================================
\note All equations are from the \textbf{LT8390 datasheet, Rev.~F}
(Applications Information, pp.~17--25) unless stated otherwise.
Values tagged \prov{} are carried over from \texttt{BuckBoostTopology.md} and
must be re-verified once the schematic is frozen.

\subsection{Input stage --- priority mux and protection}
\begin{itemize}[nosep]
  \item Priority controller: LTC4417 (or equivalent) --- fixed pin-assigned
        priority, Input~2 wired to the highest-priority pin.
  \item Per-channel UVLO: program the controller's undervoltage threshold to
        \SI{11}{\V}. \todo{insert resistor calculation from the LTC4417
        datasheet: $V_{UV}$ threshold and hysteresis formulas, show $R_{top}$, $R_{bot}$.}
  \item Fuse sizing (each input line, ahead of the mux):
        \begin{equation}
          I_{FUSE} \ge \frac{P_{OUT}}{\eta \cdot V_{IN,min}} \, k_{margin}
        \end{equation}
        with the worst-case input current evaluated at $V_{IN,min} = \SI{11}{\V}$
        and the burst power ($24\,\text{V} \times 8\,\text{A} = \SI{192}{\W}$).
        Margin factor $k_{margin} = \todo{1.5--2, to be confirmed}$.
        \todo{finalize after the efficiency estimate (\autoref{sec:efficiency})}.
\end{itemize}

\subsection{Buck-boost power stage}

\subsubsection{Switching frequency}
$R_{12} = \SI{226}{\kilo\ohm}$ (1\,\%) selects $f_{SW} = \SI{200}{\kHz}$
(datasheet Table~1). Spread spectrum ON ($\pm$15\,\% modulation) for lower EMI.
Rationale: thermal/size compromise; see \autoref{app:decisions}.

\subsubsection{Inductor ripple current}
\begin{align}
  \Delta I_{L(BUCK)}  &= \frac{V_{OUT}\,(V_{IN(MAX)} - V_{OUT})}{f_{SW}\,L\,V_{IN(MAX)}} \\[2pt]
  \Delta I_{L(BOOST)} &= \frac{V_{IN(MIN)}\,(V_{OUT} - V_{IN(MIN)})}{f_{SW}\,L\,V_{OUT}}
\end{align}
Worked worst case (\SI{11}{\V} in, \SI{24}{\V} out, $L = \SI{10}{\uH}$,
$f_{SW} = \SI{200}{\kHz}$):
\begin{equation}
  \Delta I_{L(BOOST)} = \frac{\SI{11}{\V}\,(\SI{24}{\V}-\SI{11}{\V})}
                             {\SI{200}{\kHz}\cdot\SI{10}{\uH}\cdot\SI{24}{\V}}
                       \approx \SI{2.98}{\A}
\end{equation}

\subsubsection{Inductor selection}
\begin{align}
  L_{BUCK}  &> \frac{V_{OUT}\,(V_{IN(MAX)} - V_{OUT})}{f_{SW}\,I_{OUT}\,\Delta I_L\% \, V_{IN(MAX)}} \\[2pt]
  L_{BOOST} &> \frac{V_{IN(MIN)}^2\,(V_{OUT} - V_{IN(MIN)})}{f_{SW}\,I_{OUT}\,\Delta I_L\% \, V_{OUT}^2} \\[2pt]
  \text{slope-comp.\ stability:}\quad & L > \frac{10\,V_{OUT}\,R_{SENSE}}{f_{SW}} \approx \SI{2.4}{\uH}
\end{align}
Chosen: $L_1 = \SI{10}{\uH}$, $\ge$\SI{18}{\A} saturation, ferrite, low DCR.
Height-budget tradeoff: the \SI{11}{\V}$\rightarrow$\SI{24}{\V}, \SI{8}{\A} corner
peaks at $\approx$\SI{19}{\A} --- accepted, see risk log (\autoref{app:open}).

\subsubsection{Inductor current sense and peak limit}
\begin{align}
  R_{SENSE(BUCK)}  &= \frac{2 \cdot \SI{50}{\mV}}{2\,I_{OUT} + \Delta I_{L(BUCK)}} \\[2pt]
  R_{SENSE(BOOST)} &= \frac{2 \cdot \SI{50}{\mV} \cdot V_{IN(MIN)}}
                           {2\,I_{OUT}\,V_{OUT} + \Delta I_{L(BOOST)}\,V_{IN(MIN)}}
\end{align}
Worst corner (\SI{11}{\V} in, \SI{24}{\V} out, \SI{8}{\A}):
$R_{SENSE(BOOST)} \approx \SI{2.64}{\m\ohm}$.
Chosen $R_5 = \SI{2}{\m\ohm}$ $\Rightarrow$ \textbf{24\,\% margin}
(datasheet requires final value 20--30\,\% below calculated).
Peak limit: $I_{PEAK} = \SI{50}{\mV}/R_5 = \SI{25}{\A}$.

\subsubsection{Output CC limit and monitor}
With CTRL tied to VREF ($\ge \SI{1.35}{\V}$), the ISP/ISN threshold is the full
\SI{100}{\mV} scale:
\begin{equation}
  I_{LIM} = \frac{\SI{100}{\mV}}{R_{11}}
          = \frac{\SI{100}{\mV}}{\SI{12}{\m\ohm}} = \SI{8.33}{\A}
\end{equation}
Monitor output: $V_{ISMON} = 10\,V_{(ISP-ISN)} + \SI{0.25}{\V}$.
\todo{transcription re-checked against the LT8390 datasheet p.~22.}

\subsubsection{Capacitors}
\begin{itemize}[nosep]
  \item \textbf{Input:} RMS ripple current in the buck region
        $I_{IN,RMS} \approx I_{OUT}\sqrt{D(1-D)}$, $D = V_{OUT}/V_{IN}$;
        worst case $D = 0.5$ gives $I_{OUT}/2 \approx \SI{4}{\A}$ at the 8\,A burst
        rating --- shared by \comp{C11}--\comp{C14} (4 $\times$ \SI{10}{\uF},
        \SI{50}{\V} X7R/X7S, 1210). Plus \comp{C8} = \SI{1}{\uF} at the VIN pin.
  \item \textbf{Output:} \comp{C19}--\comp{C22} (4 $\times$ \SI{10}{\uF},
        \SI{50}{\V}) + \comp{C23} (\SI{180}{\uF} polymer); \comp{C24} =
        \SI{1}{\uF} at the VOUT pin.
  \item \textbf{Bootstrap:} \comp{C10}, \comp{C18} = \SI{100}{\nF} X7R
        ($\ge$100$\times$ top-FET gate charge; datasheet range
        \SIrange{0.1}{0.47}{\uF}).
  \item \textbf{INTVCC:} \comp{C15} = \SI{4.7}{\uF} plus a second \SI{1}{\uF}
        right at the pin (datasheet layout recommendation).
\end{itemize}

\subsubsection{MOSFET selection}
\comp{Q3}--\comp{Q6} = NCEP40T11G (logic-level; gate drive is the \SI{5}{\V}
INTVCC rail). $R_{DS(ON)} \approx \SI{3.3}{\m\ohm}$ typ.\ at
$V_{GS} = \SI{4.5}{\V}$. 40\,V BVDSS on a 36\,V rail leaves $\approx$11\,\%
ringing headroom --- bench-verify, snubber if needed.

\subsection{Feedback divider and output voltage select}
\begin{equation}
  V_{OUT} = \SI{1.00}{\V} \times \frac{R_{top} + R_{bot}}{R_{bot}},
  \qquad R_{bot} = R_{17} = \SI{10.0}{\kilo\ohm}
\end{equation}
\begin{align}
  \SI{12}{\V}:&\quad R_{top} = R_{15} = \SI{110}{\kilo\ohm} \\
  \SI{24}{\V}:&\quad R_{top} = R_{15} + R_{16} = \SI{230}{\kilo\ohm}
\end{align}
\comp{J1} shunts \comp{R16}: \textbf{fitted = 12\,V, removed = 24\,V}; a missing
jumper defaults to \SI{12}{\V} (fails low, no runaway).

\subsection{Enable / UVLO}
\begin{align}
  V_{IN(turn\text{-}on)}  &= \SI{1.233}{\V}\cdot\frac{R_8+R_9}{R_9} + \SI{2.5}{\uA}\cdot R_8 \nonumber \\
  V_{IN(turn\text{-}off)} &= \SI{1.220}{\V}\cdot\frac{R_8+R_9}{R_9}
\end{align}
With $R_8 = \SI{365}{\kilo\ohm}$, $R_9 = \SI{56.2}{\kilo\ohm}$:
\textbf{$\approx$\SI{10.2}{\V} on, $\approx$\SI{9.1}{\V} off} --- deliberately
below the LTC4417's \SI{11}{\V} floor so the mux arbitrates first.

\subsection{Soft-start and compensation}
\begin{equation}
  t_{SS} = \frac{\SI{1}{\V} \cdot C_9}{\SI{12.5}{\uA}}
         = \frac{\SI{1}{\V} \cdot \SI{100}{\nF}}{\SI{12.5}{\uA}}
         \approx \SI{8}{\ms}
\end{equation}
\comp{R10} = \SI{499}{\kilo\ohm} (SS$\rightarrow$VREF) selects \textbf{latch-off}
on fault. Compensation \comp{R13} = \SI{15}{\kilo\ohm} + \comp{C25} =
\SI{15}{\nF}; bench-tune with load steps at both outputs. Constraint:
$C_9 \ge 5$--$10 \times C_{25}$ (ratio here: 6.7$\times$, OK).

\subsection{Over-temperature latch}
Trip point set by the NTC divider against the TL431's \SI{2.495}{\V} reference,
fed from a \SI{5.1}{\V} zener rail so the trip point barely moves with battery
voltage:
\begin{equation}
  R_{NTC}(trip) = R_{21}\left(\frac{V_Z}{2.495} - 1\right)
                = 1.33\text{k}\left(\frac{5.1}{2.495}-1\right)
                \approx \SI{1.39}{\kilo\ohm}
\end{equation}
$\Rightarrow$ trip $\approx$\SI{85}{\celsius} nominal, $\pm$\SI{8}{\celsius}
realistic (zener tolerance, TL431 grade, NTC B-tolerance).
\comp{TH1} (10\,k NTC) must be mounted on the hot copper, not in free air.
Latch resets only when input power is cycled.

\subsection{Indicator resistors}
Continuous worst case is 36\,V input; sizes chosen so no resistor exceeds
$\approx$50\,\% of its package rating:
\begin{table}[H]
  \centering
  \begin{tabular}{@{}llll@{}}
    \toprule
    \textbf{Ref} & \textbf{Value} & \textbf{Dissipation (worst case)} & \textbf{Package} \\
    \midrule
    \comp{R18} (green)   & \SI{4.7}{\kilo\ohm} & $\approx$\SI{0.10}{\W} & 1206 \\
    \comp{R19} (red)     & \SI{6.8}{\kilo\ohm} & $\approx$\SI{0.17}{\W} & 1206 min., 2010 preferred \\
    \comp{R29} (red)     & \SI{6.8}{\kilo\ohm} & $\approx$\SI{0.17}{\W} & 1206 min., 2010 preferred \\
    \bottomrule
  \end{tabular}
  \caption{LED series resistors. \todo{re-check currents against final LED V\textsubscript{F} bins.}}
\end{table}

\subsection{Thermal analysis}
\begin{itemize}[nosep]
  \item \textbf{INTVCC self-heating} ($\approx$\SI{1.6}{\W} inside the IC at
        36\,V, 200\,kHz, four $\approx$50\,nC gates):
        \begin{equation}
          P = V_{IN} \cdot I_{INTVCC} \approx 36 \cdot \SI{45}{\mA} \approx \SI{1.6}{\W},
          \qquad T_J = T_A + P\,\theta_{JA}
        \end{equation}
        \todo{verify $\theta_{JA}$ for the chosen layout and confirm $T_J$ at max
        ambient; fallback options: \SI{150}{\kHz}, lower-$Q_g$ FETs.}
  \item \textbf{Bottom thermal pad:} \SI{45}{\mm}$\times$\SI{40}{\mm} exposed pad,
        vias from \comp{U1} exposed pad and the power MOSFETs; contact with the
        enclosure floor. \todo{add measured case temperature once available.}
  \item \textbf{Heatsinks:} low-profile finned heatsinks on \comp{U1} and the
        switching MOSFETs where height allows.
\end{itemize}

\subsection{Efficiency budget}
\label{sec:efficiency}
\todo{fill from bench measurement or a loss budget (conduction + switching +
gate charge). One row per corner:}
\begin{table}[H]
  \centering
  \begin{tabular}{@{}lcc@{}}
    \toprule
    \textbf{Corner} & \textbf{Input} & \textbf{Output} \\
    \midrule
    Buck, 36\,V $\rightarrow$ 12\,V, \SI{2}{\A}   & \todo{\%} & \todo{\%} \\
    Buck-boost corner, 12\,V in, 12\,V out        & \todo{\%} & \todo{\%} \\
    Boost, 11\,V $\rightarrow$ 24\,V, \SI{2}{\A}  & \todo{\%} & \todo{\%} \\
    \bottomrule
  \end{tabular}
  \caption{Efficiency at nominal duty.}
\end{table}

%=============================================================================
\section{Verification \& Test Plan}
%=============================================================================
\todo{expand into a full procedure once the first board is built.}
\begin{table}[H]
  \centering
  \begin{tabularx}{\linewidth}{@{}lX@{}}
    \toprule
    \textbf{Test} & \textbf{Pass criterion} \\
    \midrule
    Input priority          & Input~2 always selected when valid; clean fallback to Input~1 \\
    UVLO per channel        & Both channels disable below \SI{11}{\V}, restart above it \\
    Reverse polarity        & No damage / no output with either input reversed \\
    Output voltage          & 12\,V and 24\,V within $\pm$1.5\,\% at nominal load \\
    Output CC limit         & Current limit at $8.33\,\text{A} \pm \text{tolerance}$ \\
    Load transient          & Stable at both outputs, no latch-off, acceptable deviation \\
    Over-temp trip          & Latch fires at target temperature; recovers only on power cycle \\
    Thermal soak            & $T_J$, case and enclosure-floor temperatures within limits at worst corner \\
    \bottomrule
  \end{tabularx}
  \caption{Verification test matrix.}
\end{table}

%=============================================================================
\section{Open Items \& Risks}
%=============================================================================
See the risk register in \autoref{app:open}.

%=============================================================================
\appendix
%=============================================================================

\section{Bill of Materials}
\label{app:bom}
\todo{Keep in sync with \texttt{SchematicBOM.md}; summarize the main parts here.}
\begin{table}[H]
  \centering
  \begin{tabularx}{\linewidth}{@{}llXl@{}}
    \toprule
    \textbf{Ref} & \textbf{Part / value} & \textbf{Purpose} & \textbf{Status} \\
    \midrule
    \comp{U1}   & LT8390 (TSSOP-28 EP)      & Buck-boost controller            & placed \\
    \comp{Q3}--\comp{Q6} & NCEP40T11G       & Power switches                   & placed \\
    \comp{L1}   & \SI{10}{\uH}, $\ge$\SI{18}{\A} sat & Main inductor             & placed \\
    \comp{R5}   & \SI{2}{\m\ohm}            & Inductor current sense           & placed \\
    \comp{R11}  & \SI{12}{\m\ohm}           & Output current sense             & placed \\
    \comp{U3} + \comp{TH1} & TL431 + NTC   & Over-temp comparator / sensor    & \todo{verify} \\
    \comp{Q7}--\comp{Q11}, \comp{D1}--\comp{D4}, \comp{J1} & \todo{parts} & Latch, LEDs, output select & to add \\
    \bottomrule
  \end{tabularx}
  \caption{Principal components (excerpt).}
\end{table}

\section{Equation Reference}
\label{app:equations}
Full derivations live in \autoref{sec:calcs}; this is a quick-reference index.
\begin{table}[H]
  \centering
  \begin{tabularx}{\linewidth}{@{}lXl@{}}
    \toprule
    \textbf{Quantity} & \textbf{Equation} & \textbf{Source} \\
    \midrule
    Inductor ripple (buck)   & $V_{OUT}(V_{IN(MAX)}-V_{OUT})/(f L V_{IN(MAX)})$   & LT8390 p.~18 \\
    Inductor ripple (boost)  & $V_{IN(MIN)}(V_{OUT}-V_{IN(MIN)})/(f L V_{OUT})$   & LT8390 p.~18 \\
    Peak current limit       & $I_{PEAK} = \SI{50}{\mV}/R_{SENSE}$                & LT8390 p.~18 \\
    Output CC limit          & $I_{LIM} = \SI{100}{\mV}/R_{11}$ (CTRL high)       & LT8390 p.~22 \\
    Output voltage           & $V_{OUT} = 1.00 \cdot (R_{top}+R_{bot})/R_{bot}$   & LT8390 p.~21 \\
    UVLO                     & $V_{IN} = 1.233\,(R_8+R_9)/R_9 + 2.5\,\text{\textmu A}\,R_8$ & LT8390 p.~21 \\
    Soft-start               & $t_{SS} = \SI{1}{\V} \cdot C_9/\SI{12.5}{\uA}$     & LT8390 p.~21 \\
    NTC trip resistance      & $R_{NTC} = R_{21}(V_Z/2.495 - 1)$                  & TL431 datasheet \\
    NCV/\todo           & \todo{}                                            & \todo{} \\
    \bottomrule
  \end{tabularx}
  \caption{Equation quick reference.}
\end{table}

\section{Design Decisions Log}
\label{app:decisions}
\begin{table}[H]
  \centering
  \begin{tabularx}{\linewidth}{@{}Xl@{}}
    \toprule
    \textbf{Decision} & \textbf{Date} \\
    \midrule
    LT8390 4-switch buck-boost: single inductor, input straddles both outputs, PGOOD + latch-off available. & 2026-09 \\
    200\,kHz switching: thermal/size compromise; spread spectrum on. & 2026-09 \\
    \SI{10}{\uH}/\SI{18}{\A}-sat inductor: height budget; accepted \SI{19}{\A} corner peak. & 2026-09 \\
    Jumper presets instead of trim pot: no drift, no knob to bump, lower height. & 2026-09 \\
    Latch pulls EN/UVLO: assured full shutdown; resets only on power cycle. & 2026-09 \\
    Two red LEDs (fault vs.\ over-temp): distinguishable states, simpler than steered diode. & 2026-09 \\
    \todo{append each new decision with rationale and date.} & \\
    \bottomrule
  \end{tabularx}
  \caption{Design decision record.}
\end{table}

\section{Open Items \& Risk Register}
\label{app:open}
\begin{table}[H]
  \centering
  \begin{tabularx}{\linewidth}{@{}Xll@{}}
    \toprule
    \textbf{Item} & \textbf{Severity} & \textbf{Status} \\
    \midrule
    Analog support blocks (\comp{Q7}--\comp{Q11}, \comp{D1}--\comp{D4}, \comp{J1}, \comp{TH1}) not yet placed & high & open \\
    No footprints assigned on any symbol & high & open \\
    Fuse ratings not finalized & medium & open \\
    INTVCC self-heating ($\approx$\SI{1.6}{\W}) --- $T_J$ margin unverified & high & open \\
    \SI{19}{\A} corner peak vs.\ \SI{18}{\A} inductor saturation & medium & accepted \\
    40\,V FETs on 36\,V rail --- ringing headroom & medium & bench-check \\
    Height budget: PCB + parts + heatsink under \SI{13.5}{\mm} & medium & open \\
    \bottomrule
  \end{tabularx}
  \caption{Open items carried from the current schematic audit.}
\end{table}

\section{Revision History}
\begin{table}[H]
  \centering
  \begin{tabular}{@{}clll@{}}
    \toprule
    \textbf{Rev} & \textbf{Date} & \textbf{Author} & \textbf{Change} \\
    \midrule
    0.1 & \today & \todo{name} & Initial template created \\
    \todo{} & \todo{} & \todo{} & \todo{} \\
    \bottomrule
  \end{tabular}
  \caption{Revision history.}
\end{table}

\section{References}
\begin{enumerate}[label={[\arabic*]},nosep]
  \item LT8390 datasheet, Rev.~F --- \texttt{Datasheets/LTC8390.pdf}
  \item LTC4417 datasheet --- \texttt{Datasheets/LTC4417.pdf}
  \item NCEP40T11G datasheet --- \texttt{Datasheets/NCEP40T11G.pdf}
  \item TL431 datasheet --- \todo{download to \texttt{Datasheets/agentDownload/}}
  \item Project specification --- \texttt{Specifications/SpecSheet.pdf}
  \item \todo{add any further references as they are used.}
\end{enumerate}

\end{document}
